\documentclass[11pt,a4paper]{article}
\usepackage[margin=2.4cm]{geometry}
\usepackage{amsmath,amssymb,amsthm,mathtools}
\usepackage{booktabs}
\usepackage{enumitem}
\usepackage{xcolor}
\usepackage{hyperref}
\usepackage{fontspec}
\usepackage{polyglossia}
\setmainlanguage{russian}
\setotherlanguage{english}
\setmainfont{Times New Roman}
\setmonofont{Consolas}[Scale=0.90]

\hypersetup{
  colorlinks=true,
  linkcolor=blue!70!black,
  citecolor=blue!70!black,
  urlcolor=blue!70!black,
  pdftitle={Целочисленная кинематика трёхвальных планетарных дифференциалов A + B = C: вронскианы Безу-Паде, понижение степени на круге и погружение в ротатор Z_4Q в Lean 4},
  pdfauthor={Каратаев Дмитрий Александрович}
}

\theoremstyle{plain}
\newtheorem{theorem}{Теорема}[section]
\newtheorem{lemma}[theorem]{Лемма}
\newtheorem{proposition}[theorem]{Предложение}
\newtheorem{corollary}[theorem]{Следствие}
\theoremstyle{definition}
\newtheorem{definition}[theorem]{Определение}
\newtheorem{remark}[theorem]{Замечание}

\title{\textbf{Целочисленная кинематика трёхвальных планетарных дифференциалов $A + B = C$:\\ вронскианы Безу--Паде, понижение степени на круге и погружение в ротатор $\mathbb{Z}_{4Q}$ в Lean~4}}
\author{Д.\,А.~Каратаев\thanks{AI-Guiders Research (\url{https://github.com/AI-Guiders}). Формальная верификация в Lean~4 и аналитическая разработка выполнены при участии ИИ-партнёра \emph{Луч} (\emph{Luch}, AI-Guiders).}}
\date{2026}

\begin{document}
\maketitle

\begin{abstract}
Всякая взаимно простая аддитивная тройка $A + B = C$ ($A = a_1 a_2,\ B = b_1 b_2,\ C = c_1 c_2$, с этажным разложением $A = \widetilde{A}_k X_k^k,\ B = \widetilde{B}_k Y_k^k,\ C = \widetilde{C}_k Z_k^k$ и оценкой весов $\widetilde{A}_k\widetilde{B}_k\widetilde{C}_k \le \operatorname{rad}(ABC)^{k-1}$) формулируется как трёхвальный планетарный дифференциал с ортогональным $3\mathrm{D}$-репером $\mathcal{F}_3 = (\mathbf{r},\mathbf{\omega},\mathbf{r}\times\mathbf{\omega})\in\mathbb{Z}^{3\times 3}$ ($\mathbf{r} = (a_1,b_1,c_1)\perp\mathbf{\omega} = (a_2,b_2,-c_2)$), плоским законом скольжения шестерён Брахмагупты $C^2 + S_{AB}^2 = (A-B)^2 + M_{AB}^2 = (a_1^2+b_1^2)(a_2^2+b_2^2)$, невырожденной матрицей моментов Ганкеля--Вандермонда $-\det\mathcal{H}_3 = \widetilde{A}\widetilde{B}\widetilde{C}(X-Y)^2(Y-Z)^2(Z-X)^2 > 0$, проекторами Лагранжа--Вандермонда зазоров валов и унимодулярной $2\times 3$-матрицей партнёров Безу $\mathcal{B}_{2\times 3}$ (\texttt{thm8}, \texttt{thm9}). Далее доказаны пять теорем в кольце $\mathbb{Z}$: обратная аннигиляция кубов Паде--Безу и квадратичный спуск $I_2$ (\texttt{thm12}); целочисленный вронскиан Мейсона--Стотерса--Безу $W^* = A B^* - B A^* \equiv X^2 Y^2 Z^2$ с определителем четвёртых степеней $-\det\mathcal{W}_{\mathrm{MS}} \equiv X^4 + Y^4 + Z^4$ (\texttt{thm13}); безусловное понижение степени остатком на круге с полным радикал-скелетом валов $R_{\mathrm{full}} = (\widetilde{A}\widetilde{B}\widetilde{C})(X Y Z)$ (\texttt{thm14}); межэтажное сцепление $k=2\leftrightarrow k=3$ на кривой Ферма--Сельмера 6-й степени рода $g(6)=10\ge 2$ (\texttt{thm15}); и стереографическое погружение Вейерштрасса в четырёхквадрантный ротатор $\mathbb{Z}_{4Q}$ ($\forall Q\in\mathbb{Z}$, включая $\mathbb{Z}_{512}$ при $Q=128$) с инвариантом четверти кольца $C\Delta_{\mathrm{rot},Q} - w I_{\mathrm{rot},Q} \equiv Q$ (\texttt{thm16}). Все теоремы верифицированы в \texttt{AbcPlanetaryWronskian.lean} (\texttt{0 sorryAx}, \texttt{0 Classical.choice}).

\medskip\noindent
\textbf{Ключевые слова:} аддитивные диофантовы тройки, $abc$-тройки, вронскиан Мейсона--Стотерса, аннигиляция Паде--Безу, кривые Ферма--Сельмера, стереографический гомеоморфизм Вейерштрасса, Lean~4.
\end{abstract}

\section{Репер планетарного дифференциала \texorpdfstring{$\mathcal{F}_3$}{F3}, матрица Ганкеля--Вандермонда \texorpdfstring{$\mathcal{H}_3$}{H3} и трёхканальная матрица Паде \texorpdfstring{$\mathcal{P}_3$}{P3}}
\label{sec:planetary-frame}

Следуя терминологии теории моделей А.\,И.~Мальцева (\emph{«Алгебраические системы»}~\cite{Malcev1970}, где \emph{носителем} / \emph{carrier} называется основное множество алгебраического сорта), мы рассматриваем систему как двусортную алгебраическую структуру $\langle\mathcal{M},\mathcal{T};\,\pi_{M\to T},\,\iota_{T\to M}\rangle$, связывающую \textbf{дискретный $M$-носитель} ($\mathbb{Z}_{4Q}$ с $L^1$-проекциями $\mathcal{C}_{4Q},\mathcal{S}_{4Q}$) и \textbf{непрерывный евклидов $T$-носитель} ($L^2$-окружность $X^2+Y^2=R^2$ и прямую $\mathbb{R}$).

\begin{theorem}[Окаймлённая $3\times 3$-матрица $\mathcal{M}_3$, ортогональный репер дифференциала $\mathcal{F}_3$ и матрица Ганкеля--Вандермонда $\mathcal{H}_3$, \texttt{thm8}]
\label{thm:main-8}
Для всякой аддитивной тройки $A + B = C$ ($A = a_1 a_2,\ B = b_1 b_2,\ C = c_1 c_2$) матрица $\mathcal{M}_3 = \begin{pmatrix} c_2 & a_2 & b_2 \\ a_1 & c_1 & 0 \\ b_1 & 0 & c_1 \end{pmatrix}$ имеет $\det\mathcal{M}_3 = c_1(C-A-B) = 0$ и союзную матрицу ранга~$1$ $(\mathcal{M}_3^{\mathrm{adj}})^2 = (c_1^2+C)\mathcal{M}_3^{\mathrm{adj}}$. Вектор радиусов шестерён $\mathbf{r} = (a_1,b_1,c_1)$ и вектор чисел оборотов $\mathbf{\omega} = (a_2,b_2,-c_2)$ ортогональны ($\mathbf{r}\cdot\mathbf{\omega} = 0$), $\det\mathcal{F}_3 = M_{BC}^2 + M_{CA}^2 + S_{AB}^2 \equiv \|\mathbf{r}\|_2^2\|\mathbf{\omega}\|_2^2$, $C^2 + S_{AB}^2 = (A-B)^2 + M_{AB}^2 \equiv (a_1^2+b_1^2)(a_2^2+b_2^2)$, а на этаже $k=3$ ($S_m = \widetilde{A}X^m + \widetilde{B}Y^m - \widetilde{C}Z^m$, $S_3 = 0$) матрица Ганкеля $\mathcal{H}_3$ имеет $-\det\mathcal{H}_3 \equiv \widetilde{A}\widetilde{B}\widetilde{C}(X-Y)^2(Y-Z)^2(Z-X)^2$ и проекторы Лагранжа--Вандермонда зазоров валов $\widetilde{A}X(X-Y)(X-Z) \equiv S_1 YZ - S_2(Y+Z)$, $\widetilde{B}Y(Y-X)(Y-Z) \equiv S_1 ZX - S_2(Z+X)$ и $\widetilde{C}Z(Z-X)(Z-Y) \equiv S_2(X+Y) - S_1 XY$.
\end{theorem}

\begin{theorem}[Унимодулярная $2\times 3$-матрица партнёров Безу $\mathcal{B}_{2\times 3}$ и трёхканальная матрица Паде $\mathcal{P}_3(A,B,C)$, \texttt{thm9}]
\label{thm:main-9}
Для любых $A + B = C$ и пары Безу $A v - B u = 1$ ($w = u + v$) все три минора $2\times 2$ матрицы $\mathcal{B}_{2\times 3} = \begin{pmatrix} A & B & C \\ u & v & w \end{pmatrix}$ равны $+1$ ($Av-Bu = Aw-Cu = Bw-Cv = 1$), а трёхканальная кубическая матрица Паде $\mathcal{P}_3(A,B,C)$ имеет ядро $\mathcal{P}_3(A,B,C)(A,B,-C)^T = \mathbf{0}$ и след $\operatorname{Tr}\mathcal{P}_3(A,B,C) = A C(A+C)$.
\end{theorem}

\section{Вронскиан Мейсона--Стотерса, понижение степени на круге, замок рода~10 и погружение в \texorpdfstring{$\mathbb{Z}_{4Q}$}{Z\_4Q}}
\label{sec:universal-theorems}

\begin{theorem}[Обратная аннигиляция кубов Паде--Безу, исключение медианты Штерна--Броко и квадратичный спуск $I_2$, \texttt{thm12}]
\label{thm:reverse-pade-thm12}
Для $A = \widetilde{C}Z^3 - \widetilde{B}Y^3$, $I_0 = Yw - Zv$ и $I_1 = Yw(\widetilde{C}w^3+2\widetilde{B}v^3) - Zv(2\widetilde{C}w^3+\widetilde{B}v^3)$:
\begin{equation}
  \label{eq:reverse-pade-bezout}
  Z^3 I_1 \;+\; A\,w^3(2Zv - Yw) \;\equiv\; \widetilde{B}\,I_0^3(Yw + Zv),\qquad
  Y^3 I_1 \;+\; A\,v^3(2Yw - Zv) \;\equiv\; \widetilde{C}\,I_0^3(Yw + Zv),
\end{equation}
а между двумя родителями Штерна--Броко $(v,w)$ и $(v_m,w_m) = (Y-v,Z-w)$ медианты Фарея $(Y,Z)$ ($p = v_m w,\ q = v w_m,\ K_+ = w^3(2Zv-Yw),\ K_- = w_m^3(2Zv_m-Yw_m)$) исключение $A$ сокращает $Z^3$ тождественно: $I_1 K_- - I_{1m} K_+ \equiv \widetilde{B} I_0^3\bigl((vw + v_m w_m)(p+q) - (p^2+q^2)\bigr)$. Инвариант квадратичного спуска $I_2 := \widetilde{B}Yv^2(3Yw-Zv) - \widetilde{C}Zw^2(3Zv-Yw)$ исключает родителя $(v,w)$ из правой части ($Z^2 I_2 + A w^2(3Zv-Yw) \equiv \widetilde{B}Y I_0^3$, $Y^2 I_2 + A v^2(3Yw-Zv) \equiv \widetilde{C}Z I_0^3$) и образует с $I_1$ не зависящую от $A$ систему Вронскиана $2\times 2$ определителя $3 I_0^3$.
\end{theorem}

\begin{theorem}[Кубическая аннигиляция по счётной башне аппроксимантов Паде, исключение $Z^3$ и выход из нулевой ячейки между соседними этажами $\mathcal{R}\in\mathrm{List}\ \mathbb{Z}$, \texttt{thm12b\_abc\_countable\_pade\_tower\_cubic\_annihilation\_and\_adjacent\_floor\_escape}]
\label{thm:countable-pade-tower-thm12b}
Для каждого этажа $\mathcal{R}\in\mathrm{List}\ \mathbb{Z}$ счётной башни аппроксимантов Паде $(P_{\mathcal{R}},Q_{\mathcal{R}}) = (\texttt{padeTowerP}\ \mathcal{R}\ U\ V,\ \texttt{padeTowerQ}\ \mathcal{R}\ U\ V)$, построенной над $U = \widetilde{C}Z^3$, $V = \widetilde{B}Y^3$, $A = U - V$ ($A + B = C$), и любой пары валов $(w,v)$ с $a = v Z$, $b = w Y$, $A' = \widetilde{C}w^3 - \widetilde{B}v^3$ и инвариантом этажа $I_{\mathcal{R}} = a P_{\mathcal{R}} - b Q_{\mathcal{R}}$:
\begin{enumerate}[label=\textbf{(\roman*)},leftmargin=*,noitemsep]
\item $P_{\mathcal{R}} - Q_{\mathcal{R}} = \Pi_-(\mathcal{R})\,A$, $P_{\mathcal{R}} + Q_{\mathcal{R}} = \Pi_+(\mathcal{R})\,(U + V)$, а на базовом этаже $\mathcal{R} = [3]$ выполнено $U P_{[3]}^3 - V Q_{[3]}^3 = A^3(U + V)$;
\item При $u_Z = a P_{\mathcal{R}}$, $v_Y = b Q_{\mathcal{R}}$ и $\Phi_3(u_Z,v_Y) = u_Z^2 + u_Z v_Y + v_Y^2$ для всех $Z^3 \ge 1$ множитель $Z^3$ тождественно сокращается на каждом этаже $\mathcal{R}$:
  \begin{equation}
    \label{eq:abc-arbitrary-floor-z3-cancel}
    \widetilde{C}\,I_{\mathcal{R}}\,\Phi_3(u_Z, v_Y) \;=\; v^3\bigl(U P_{\mathcal{R}}^3 - V Q_{\mathcal{R}}^3\bigr) \;-\; A'\,Y^3 Q_{\mathcal{R}}^3;
  \end{equation}
\item Между соседними этажами $\mathcal{R}$ и $\mathcal{R}' = \mathcal{R}\mathbin{++}[k+1]$ выполнено $I_{\mathcal{R}'} - (k+1)I_{\mathcal{R}} = -k(a+b)\Pi_-(\mathcal{R})A$, откуда при $k(a+b)\Pi_-(\mathcal{R})A \ne 0$ имеем $\neg(I_{\mathcal{R}} = 0 \land I_{\mathcal{R}'} = 0)$, $|I_{\mathcal{R}}| + |I_{\mathcal{R}'}| \ge 1$ и $I_{\mathcal{R}} = 0 \implies I_{\mathcal{R}'} = -k(a+b)\Pi_-(\mathcal{R})A \ne 0$.
\end{enumerate}
\end{theorem}

\begin{theorem}[Целочисленный кубический вронскиан Мейсона--Стотерса--Безу и определитель четвёртых степеней, \texttt{thm13}]
\label{thm:mason-stothers-thm13}
При $\widetilde{A}X^3 + \widetilde{B}Y^3 = \widetilde{C}Z^3$ и $\widetilde{A}Xv_0 - \widetilde{B}Yu_0 = 1$ подстановка $\alpha = u_0 Z^2 + t\widetilde{A}X$, $\beta = v_0 Z^2 + t\widetilde{B}Y$, $\gamma = u_0 X^2 + v_0 Y^2 + t\widetilde{C}Z$ даёт $A^*+B^*=C^*$ ($A^*=\alpha X^2,\ B^*=\beta Y^2,\ C^*=\gamma Z^2$) с минорами квадратов валов $(\widetilde{A}X\beta-\widetilde{B}Y\alpha,\ \widetilde{C}Z\alpha-\widetilde{A}X\gamma,\ \widetilde{C}Z\beta-\widetilde{B}Y\gamma) = (Z^2,Y^2,X^2)$ и точным вронскианом
\begin{equation}
  \label{eq:ms-wronskian}
  W^* \;:=\; A B^* - B A^* \;\equiv\; X^2 Y^2 Z^2,\qquad
  -\det\mathcal{W}_{\mathrm{MS}} \;\equiv\; X^4 + Y^4 + Z^4.
\end{equation}
\end{theorem}

\begin{theorem}[Понижение степени остатком на круге с полным радикал-скелетом валов и кубический замок Крамера--Безу $3\times 3$, \texttt{thm14}]
\label{thm:circle-remainder-thm14}
Для полного радикал-скелета с валами $R_{\mathrm{full}} := (\widetilde{A}\widetilde{B}\widetilde{C})(X Y Z)$ выполнено объёмное тождество $A B C \equiv R_{\mathrm{full}} W^*$ и круговое разложение
\begin{equation}
  \label{eq:circle-remainder-rfull}
  (\widetilde{B}Y)\,C \;+\; (\widetilde{B}^2\widetilde{C})(Y^2 Z)\,\alpha \;\equiv\; R_{\mathrm{full}}\,\beta,\qquad
  (\widetilde{A}X)\,\beta \;=\; Z^2 + (\widetilde{B}Y)\,\alpha.
\end{equation}
Выбор неотрицательного остатка $0\le\alpha\le\widetilde{A}X-1$ на круге длины $\widetilde{A}X$ даёт $(\widetilde{A}X)\beta \le Z^2 + (\widetilde{B}Y)(\widetilde{A}X-1)$ и $(\widetilde{B}Y)C \le R_{\mathrm{full}}\beta$ (откуда $C \le R_{\mathrm{full}}^2$ при $\beta\le R_{\mathrm{full}}$ и $C \le R_{\mathrm{full}}^3$ для всех положительных валов и весов). Для $3\times 3$ кубической системы Крамера--Безу $U_0 Y^3 - V_0 X^3 = 1$ и $V_1 Z^3 - U_1 Y^3 = 1$ определитель $\Delta_{\mathrm{cube}} := \widetilde{A}U_0 V_1 + \widetilde{B}V_0 V_1 - \widetilde{C}U_1 V_0 \in [1, \widetilde{A}+\widetilde{C}]$ удовлетворяет $\Delta_{\mathrm{cube}}(X^3, Y^3, Z^3) = \bigl(\widetilde{C}(U_0+U_1)-\widetilde{B}V_1,\ \widetilde{A}V_1+\widetilde{C}V_0,\ \widetilde{A}(U_0+U_1)+\widetilde{B}V_0\bigr)$.
\end{theorem}

\begin{theorem}[Трихотомия дефектов нижних этажей Лагранжа--Вандермонда $(S_1,S_2,S_3=0)$, трёхствольное замыкание миноров $(Z^2,Y^2,X^2)$ и совместная двусторонняя вилка Безу для $(\alpha,\beta,\gamma)$ в главной ячейке, \texttt{thm14b}, \texttt{thm14c}]
\label{thm:floor-defect-trichotomy-thm14b}
Для любого радикально-этажного разложения $A = \widetilde{A}X^3,\ B = \widetilde{B}Y^3,\ C = \widetilde{C}Z^3$ с замыканием третьего этажа $S_3 = A + B - C = 0$ и дефектами нижних этажей $S_1 = \widetilde{A}X + \widetilde{B}Y - \widetilde{C}Z$, $S_2 = \widetilde{A}X^2 + \widetilde{B}Y^2 - \widetilde{C}Z^2$:
\begin{enumerate}[label=\textbf{(\roman*)},leftmargin=*,noitemsep]
\item \textbf{Проекторы Лагранжа--Вандермонда для полных высот при $S_3 = 0$:}
  \[
    C(Z-X)(Z-Y) \;=\; Z^2\bigl(S_2(X+Y) - S_1 XY\bigr), \quad
    A(X-Y)(X-Z) \;=\; X^2\bigl(S_1 YZ - S_2(Y+Z)\bigr), \quad
    B(Y-X)(Y-Z) \;=\; Y^2\bigl(S_1 ZX - S_2(Z+X)\bigr);
  \]
\item \textbf{Невырожденность нижних этажей и знаковая трихотомия $(S_1,S_2)$:}
  при $\widetilde{C}\ge 1,\ X\ge 1,\ Y\ge 1,\ X<Z,\ Y<Z$ выполнено $S_2(X+Y) - S_1 XY = \widetilde{C}Z(Z-X)(Z-Y) \ge 1$, откуда $\neg(S_1=0\land S_2=0)$ и $|S_1|+|S_2|\ge 1$:
  (а) при $S_1 \ge 0$ имеем $\widetilde{C}Z \le \widetilde{A}X + \widetilde{B}Y$ и $C \le (\widetilde{A}X + \widetilde{B}Y)Z^2$;
  (б) при $S_2 \le 0$ имеем $S_1 \le -1$ и $(-S_1)XY \ge 1$;
  (в) при $S_1 \le -1$ и $S_2 \ge 1$ имеем $\widetilde{A}X + \widetilde{B}Y + 1 \le \widetilde{C}Z$ и $\widetilde{C}Z^2 + 1 \le \widetilde{A}X^2 + \widetilde{B}Y^2$;
\item \textbf{Трёхствольное замыкание миноров $(Z^2,Y^2,X^2)$ и совместная двусторонняя вилка в главной ячейке ($0 \le \alpha \le \widetilde{A}X - 1$):}
  когда тройка $(\alpha,\beta,\gamma)$ удовлетворяет всем трём минорным законам $(\widetilde{A}X)\beta - (\widetilde{B}Y)\alpha = Z^2$, $(\widetilde{A}X)\gamma - (\widetilde{C}Z)\alpha = Y^2$, $(\widetilde{C}Z)\beta - (\widetilde{B}Y)\gamma = X^2$ и балансу производных $X^2\alpha + Y^2\beta = Z^2\gamma$, система кругового коллапса объёмов над $R_{\mathrm{full}} = \widetilde{A}\widetilde{B}\widetilde{C}XYZ$ имеет вид $(\widetilde{B}Y)C + (\widetilde{B}^2\widetilde{C}Y^2 Z)\alpha = R_{\mathrm{full}}\beta$, $(\widetilde{C}Z)B + (\widetilde{B}\widetilde{C}^2 Y Z^2)\alpha = R_{\mathrm{full}}\gamma$ и $(\widetilde{B}Y)(\widetilde{C}Z)A = R_{\mathrm{full}}((\widetilde{C}Z)\beta - (\widetilde{B}Y)\gamma)$, а в главной фундаментальной ячейке $0 \le \alpha \le \widetilde{A}X - 1$ ($\widetilde{A}X,\widetilde{B}Y,\widetilde{C}Z,Y,Z \ge 1$) выполнены $\beta \ge 1$, $\gamma \ge 1$, $Y^2\beta \le Z^2\gamma \le Y^2\beta + A - X^2$ и
  \begin{equation}
    \label{eq:bilateral-bezout-height-bracket}
    \begin{aligned}
      Z^2 &\le (\widetilde{A}X)\beta \le Z^2 + (\widetilde{B}Y)(\widetilde{A}X - 1), &\quad
      R_{\mathrm{full}}\beta - (\widetilde{B}^2\widetilde{C}Y^2 Z)(\widetilde{A}X - 1) &\le (\widetilde{B}Y)C \le R_{\mathrm{full}}\beta, \\
      Y^2 &\le (\widetilde{A}X)\gamma \le Y^2 + (\widetilde{C}Z)(\widetilde{A}X - 1), &\quad
      R_{\mathrm{full}}\gamma - (\widetilde{B}\widetilde{C}^2 Y Z^2)(\widetilde{A}X - 1) &\le (\widetilde{C}Z)B \le R_{\mathrm{full}}\gamma.
    \end{aligned}
  \end{equation}
\end{enumerate}
\end{theorem}

\begin{theorem}[Межэтажное сцепление $k=2 \leftrightarrow k=3$ на кривой рода~10 и трихотомия эйлеровой характеристики, \texttt{thm15}]
\label{thm:genus-10-thm15}
Эйлерова характеристика плоской кривой степени $k$ удовлетворяет тождеству $\chi(\mathcal{C}_k) = 2 - (k-1)(k-2) \equiv k(3-k)$ (обращаясь в нуль при $k = 3$ в соответствии с разложением Пуанкаре--Хегора 3-сферы $S^3 = \mathcal{C}_{13}^{(+)}\cup_{O_2(3)}\mathcal{C}_{13}^{(-)}$, где $\chi(\mathcal{C}_{13}) = +1,\ \chi(O_2(3)) = +2,\ \chi(S^3) = 0$). Сцепление сферического этажа $k=2$ ($A = A_2 X_2^2,\ B = B_2 Y_2^2,\ C = C_2 Z_2^2$, $\chi=+2$) и торового этажа $k=3$ ($A = A_3 X_3^3,\ B = B_3 Y_3^3,\ C = C_3 Z_3^3$, $\chi=0$) в силу $3-2=+1$ поднимает $(U_6,V_6,W_6) = (X_2 Y_3 Z_3, Y_2 X_3 Z_3, Z_2 X_3 Y_3)$ на кривую 6-й степени рода $g(6) = 10 \ge 2$ ($\chi(\mathcal{C}_6) = -18 < 0$):
\begin{equation}
  \label{eq:genus-10-curve}
  A_2^3 B_3^2 C_3^2\,U_6^6 \;+\; B_2^3 A_3^2 C_3^2\,V_6^6 \;=\; C_2^3 A_3^2 B_3^2\,W_6^6.
\end{equation}
\end{theorem}

\begin{theorem}[Стереографическое погружение Вейерштрасса в четырёхквадрантный ротатор $\mathbb{Z}_{4Q}$, \texttt{thm16}]
\label{thm:main-12-16}
Для любой тройки $A+B=C\in\mathbb{Z}^3$ с партнёром Безу $A v - B u = 1$ ($u+v=w$) подъём Вейерштрасса $X_T = A^2-B^2,\ Y_T = 2AB,\ R_T = A^2+B^2$ удовлетворяет $X_T^2+Y_T^2=R_T^2$, $R_T+Y_T=C^2$ и $(C^2+X_T,C^2-X_T)=(2CA,2CB)$, а для любого четверть-модуля $Q\in\mathbb{Z}$ на кольце $\mathbb{Z}_{4Q}$ (включая $\mathbb{Z}_{512}$ при $Q = 128$) и любой ячейки ротатора $\varphi_0\in\mathbb{Z}$ ($\mathcal{C}_{4Q}(\varphi_0)=Q-\varphi_0,\ \mathcal{S}_{4Q}(\varphi_0)=\varphi_0,\ I_{\mathrm{rot},Q} = B\mathcal{C}_{4Q}(\varphi_0)-A\mathcal{S}_{4Q}(\varphi_0),\ \Delta_{\mathrm{rot},Q} = \mathcal{C}_{4Q}(\varphi_0)v-\mathcal{S}_{4Q}(\varphi_0)u$) выполняется точный инвариант четверти кольца:
\begin{equation}
  \label{eq:quarter-ring-invariant}
  Q\,A \;=\; C\,\mathcal{C}_{4Q}(\varphi_0) - I_{\mathrm{rot},Q},\qquad
  Q\,B \;=\; C\,\mathcal{S}_{4Q}(\varphi_0) + I_{\mathrm{rot},Q},\qquad
  C\cdot\Delta_{\mathrm{rot},Q} \;-\; w\cdot I_{\mathrm{rot},Q} \;\equiv\; Q.
\end{equation}
\end{theorem}

\section{Формальная верификация в Lean~4}
Формальное доказательство Теорем~\ref{thm:main-8}--\ref{thm:main-12-16} (\texttt{thm8}, \texttt{thm9}, \texttt{thm12}, \texttt{thm12b\_abc\_countable\_pade\_tower\_cubic\_annihilation\_and\_adjacent\_floor\_escape}, \texttt{thm13}, \texttt{thm14}, \texttt{thm14b\_abc\_hankel\_vandermonde\_floor\_defect\_trichotomy\_and\_bilateral\_bezout\_bracket}, \texttt{thm14c\_abc\_three\_shaft\_minor\_closure\_and\_joint\_beta\_gamma\_brackets}, \texttt{thm15}, \texttt{thm16}) содержится в модуле \texttt{AbcPlanetaryWronskian.lean} в репозитории \url{https://github.com/AI-Guiders/scientific-papers} (\texttt{0 sorryAx}, \texttt{0 Classical.choice}).

\begin{thebibliography}{9}
\bibitem{Malcev1970}
А.\,И.~Мальцев, \emph{Алгебраические системы}, М.: Наука, 1970 (англ.\ пер.: A.\,I.~Mal'cev, \emph{Algebraic Systems}, Springer-Verlag, 1973).
\bibitem{Mason1984}
R.\,C.~Mason, \emph{Diophantine Equations over Function Fields}, Cambridge University Press, 1984.
\bibitem{DarmonGranville1995}
H.~Darmon and A.~Granville, ``On the equations $z^m = F(x,y)$ and $Ax^p + By^q = Cz^r$,'' \emph{Bull.~London Math.~Soc.} \textbf{27}(6) (1995), 513--543.
\bibitem{Moura2021}
L.~de~Moura and S.~Ullrich, ``The Lean~4 theorem prover and programming language,'' in \emph{CADE~28}, LNCS vol.~12699, Springer, 2021, pp.~625--635.
\end{thebibliography}

\end{document}
